\subsection{Полулинейные уравнения}

\stepcounter{equation}
\begin{equation}\tag{1}\label{eq:semilin}
  \sum_{i} a_i(x)\, u_{x_i} \equiv \bar a \cdot \nabla u = f(x, u) .
\end{equation}

\begin{warn}
Номер параграфа не проставлен ни на доске, ни в тетради: заголовок идёт сразу
после контрпримера с касанием. Здесь он взят следующим по порядку. В тетради
он записан как «Полулинейное ур-е ЧП 1-го \emph{рода}» — по смыслу речь, как и
везде, о первом \emph{порядке}. Нумерация формул в этом параграфе начинается
заново.
\end{warn}

\begin{idea}
Отличие от~\eqref{eq:fo-homog} только в правой части, причём $f$ зависит от $u$,
но \emph{не} от производных — поэтому «полулинейное». Левая часть по-прежнему
линейна по $\nabla u$, значит характеристики те же самые: они определяются одним
лишь полем $\bar a(x)$ и от $f$ не зависят. Меняется лишь то, что вдоль
характеристики $u$ теперь не постоянна, а подчиняется ОДУ.
\end{idea}

\subsubsection*{Задача Коши на поверхности}

Пусть задана поверхность размерности $n-1$:
\[
  S = \Big\{\, \vec\xi(\vec\tau) \equiv
      \big(\xi_1(\vec\tau),\dots,\xi_n(\vec\tau)\big), \quad
      \vec\tau = (\tau_1,\dots,\tau_{n-1}) \,\Big\} ,
\]
и начальное условие
\[
  u|_S = \varphi(\vec\tau) .
\]

\begin{definition}
\textbf{Условие трансверсальности}: поле направлений (то есть характеристики)
нигде не касается $S$.
\end{definition}

\inputfig{figures/fig-surface-transversal}

\subsubsection*{Замена переменных вдоль характеристик}

Выпустим из каждой точки поверхности характеристику:
\[
  \begin{cases}
    \dfrac{d\vec x}{ds} = \bar a(\vec x), \\[8pt]
    \vec x(0) = \vec\xi(\vec\tau) .
  \end{cases}
\]
Это задаёт замену переменных $(s, \vec\tau) \mapsto \vec x$, при которой сама
поверхность получает уравнение
\[
  S = \{\, s = 0 \,\} .
\]

Якобиан этой замены равен
\[
  \begin{pmatrix}
    x_{1s} & x_{1\tau_1} & \dots & x_{1\tau_{n-1}} \\
    \vdots & \vdots      &       & \vdots \\
    x_{ns} & x_{n\tau_1} & \dots & x_{n\tau_{n-1}}
  \end{pmatrix}
  \Bigg|_{s=0}
  = \Big(\, \bar a\big(\vec\xi(\vec\tau)\big) \ \Big|\
    \vec\xi_{\tau_1}(\vec\tau) \ \Big|\ \dots \ \Big|\
    \vec\xi_{\tau_{n-1}}(\vec\tau) \,\Big) \ \ne\ 0 ,
\]
и отличие определителя от нуля — это в точности условие трансверсальности.

\begin{idea}[title={Что означает этот якобиан}]
При $s = 0$ первый столбец равен
\[
  \dd{\vec x}{s}\bigg|_{s=0} = \bar a\big(\vec\xi(\vec\tau)\big) ,
\]
то есть направляющему вектору характеристики, а остальные $n-1$ столбцов —
это $\partial \vec\xi / \partial \tau_j$, то есть базис касательной плоскости к
$S$. Определитель отличен от нуля тогда и только тогда, когда $\bar a$ не лежит
в касательной плоскости, — а это и есть условие трансверсальности.

\medskip
Значит по теореме об обратной функции замена $(s,\vec\tau) \leftrightarrow
\vec x$ локально взаимно однозначна и гладко обратима в окрестности $S$. Именно
здесь трансверсальность и нужна; без неё характеристики «ложатся» на $S$,
замена вырождается, и получается ситуация вроде контрпримера с $y = x^3$
из \S1.
\end{idea}

\begin{warn}
В тетради справа от матрицы стояло просто «$= \bar a(\vec\xi)$» — отсюда и
возникал вопрос. Полная запись есть на другой странице тетради и приведена
выше: первый столбец равен $\bar a(\vec\xi)$, остальные — производные
$\vec\xi$ по $\tau_j$, а всё вместе отлично от нуля по условию
трансверсальности.
\end{warn}

\subsubsection*{Сведение к ОДУ}

Итак, обратное отображение определено, и переменные переходят друг в друга:
\[
  \vec x \ \longleftrightarrow\ (s, \vec\tau) .
\]
Сделаем замену
\[
  u(\vec x) = v(s, \vec\tau) .
\]
По соотношению~\eqref{eq:char-deriv} производная по $s$ — это в точности
$\bar a \cdot \nabla u$, поэтому задача~\eqref{eq:semilin} превращается в
\[
  \begin{cases}
    \dfrac{\partial v}{\partial s}
      \overset{\eqref{eq:char-deriv}}{=} \bar a \cdot \nabla u
      = f(x, u) = \widetilde f(s, \vec\tau, v), \\[10pt]
    v|_{s=0} = \varphi(\vec\tau) .
  \end{cases}
\]
А это обыкновенное дифференциальное уравнение — решать умеем.

\begin{idea}
Это обыкновенное дифференциальное уравнение по $s$, в котором $\vec\tau$ —
просто набор параметров. То есть метод характеристик целиком: уравнение в
частных производных распадается на семейство задач Коши для ОДУ, по одной на
каждую характеристику, а начальные данные для них считываются с поверхности
$S$.

\medskip
Линейный однородный случай \S1 — частный: там $f \equiv 0$, уравнение
принимает вид $\partial v / \partial s = 0$, то есть $v$ постоянна вдоль
характеристики, и решение сводится к переносу данных с $S$.
\end{idea}
